\documentclass[11pt]{article}
\usepackage[margin=1in]{geometry}
\usepackage{booktabs}
\usepackage{array}
\usepackage{amsmath}
\usepackage{hyperref}
\hypersetup{colorlinks=true, linkcolor=blue, urlcolor=blue, citecolor=blue}

\title{Fairness Metrics Toolkit \\ \large A User Guide}
\author{Pedro Salas-Rojo}
\date{}

\begin{document}
\maketitle

\section{What this is}

This document is the user guide for the \textbf{Fairness Metrics Toolkit}, a
self-contained R implementation of 13 binary (two-group) fairness metrics and
4 multidimensional (N-subgroup) fairness metrics used to audit whether an
algorithm's predictions treat a protected group fairly relative to a
reference group.

The toolkit is deliberately \textbf{metrics only}: it measures fairness, it
does not implement or discuss methods to correct an unfair algorithm. It
requires no packages beyond base R (\texttt{stats}), and ships with a
synthetic example dataset so every function can be tried immediately.

\subsection{Notation used throughout}

\begin{itemize}
  \item $G \in \{0,1\}$: protected attribute. $G=0$ is the disadvantaged
        group, $G=1$ the advantaged/reference group.
  \item $Y \in \{0,1\}$: true outcome.
  \item $\hat Y \in \{0,1\}$: the algorithm's decision (denoted \texttt{Yhat}
        in the code).
  \item $\mathrm{TPR} = P(\hat Y=1 \mid Y=1)$, $\mathrm{FPR} = P(\hat Y=1 \mid Y=0)$,
        $\mathrm{TNR} = P(\hat Y=0 \mid Y=0)$, $\mathrm{FNR} = P(\hat Y=0 \mid Y=1)$.
  \item $\mathrm{PPV} = P(Y=1 \mid \hat Y=1)$ (precision), $\mathrm{NPV} = P(Y=0 \mid \hat Y=0)$.
  \item $\mathrm{Acc} = P(\hat Y = Y)$.
  \item $zg_m$: subgroup $m$ in an integer-coded partition \texttt{ZG} used by
        the 4 multidimensional metrics (any number of subgroups).
  \item $\varepsilon$ (\texttt{eps}): discrimination-aversion parameter for the
        multidimensional metrics. Larger $\varepsilon$ relaxes the fairness
        threshold; $\varepsilon = 0$ is the strictest case.
\end{itemize}

\section{Repository structure and quick start}

\begin{verbatim}
fairness_metrics/
|-- README.md
|-- fairness_functions.R                 <- every metric function
|-- 01_load_data.R                       <- loads + validates the data
|-- 02_compute_metrics_one_by_one.R      <- each metric, explained individually
|-- 03_compute_all_metrics.R             <- every metric at once + summary CSV
|-- data/mock_data.csv                   <- ships pre-built
`-- docs/fairness_metrics_guide.pdf      <- this document
\end{verbatim}

\texttt{data/mock\_data.csv} ships pre-built, so no generation step is needed.
From the repository root, run the scripts in order:

\begin{verbatim}
Rscript 02_compute_metrics_one_by_one.R
Rscript 03_compute_all_metrics.R
\end{verbatim}

Running \texttt{03\_compute\_all\_metrics.R} writes a summary to
\texttt{output/fairness\_metrics\_summary.csv}, created automatically on first
run.

To apply the toolkit to your own data, build a data frame with the columns
described in Table~\ref{tab:datareq} and call \texttt{compute\_binary\_metrics()}
and/or \texttt{compute\_all\_multidimensional()} directly -- nothing else in
the repository needs to change.

\begin{table}[h]
\centering
\begin{tabular}{lll}
\toprule
Column & Used by & Requirement \\
\midrule
\texttt{G}    & 13 binary metrics & exactly 2 values (any coding) \\
\texttt{Y}    & 13 binary metrics & true outcome, binary (0/1) \\
\texttt{Yhat} & 13 binary metrics & algorithm's decision, binary (0/1) \\
\texttt{A}    & Conditional Equal Opportunity (optional) & any categorical variable \\
\texttt{ZG}   & 4 multidimensional metrics & integer subgroup code $0,1,2,\dots$ \\
\bottomrule
\end{tabular}
\caption{Data requirements to reuse the toolkit on new data.}
\label{tab:datareq}
\end{table}

\section{The worked example: a loan-approval algorithm}

\texttt{data/mock\_data.csv} was built by simulating 5{,}000 loan applicants.
Each has covariates (income, credit score, employment years) that determine
whether they would truly repay a loan ($Y$), and an algorithm predicts a
repayment probability (\texttt{score}) and approves the loan if that
probability clears 0.5 ($\hat Y$). The simulation builds in two sources of
disparity across the protected attribute $G$ (0 = disadvantaged, 1 =
advantaged):

\begin{enumerate}
  \item a structural gap in the applicant covariates themselves, reflecting a
        real-world disparity upstream of any algorithm; and
  \item a direct algorithmic penalty against $G=0$ on top of that -- the kind
        of bias a fairness audit is meant to catch.
\end{enumerate}

A second attribute, \texttt{region}, is crossed with $G$ into four subgroups
(\texttt{ZG} $= 2\cdot G + \texttt{region} \in \{0,1,2,3\}$) used by the
multidimensional metrics, and \texttt{age\_band} is used to illustrate
Conditional Equal Opportunity. No real data is used anywhere in this
repository. With seed 123, the resulting sample has:

\begin{center}
\begin{tabular}{lcc}
\toprule
 & $G=0$ (disadvantaged) & $G=1$ (advantaged) \\
\midrule
$N$                    & 2{,}255 & 2{,}745 \\
True repayment rate, $P(Y=1)$   & 0.356 & 0.564 \\
Approval rate, $P(\hat Y=1)$    & 0.251 & 0.616 \\
\bottomrule
\end{tabular}
\end{center}

All numbers reported in the rest of this guide are computed directly from
this dataset by \texttt{03\_compute\_all\_metrics.R}.

\section{The 13 binary fairness metrics}

Metrics ending in \texttt{\_diff} have expected value $0$ under parity;
metrics ending in \texttt{\_ratio} have expected value $1$. Every metric
compares $G=0$ against $G=1$.

\begin{table}[h]
\centering
\small
\begin{tabular}{p{0.3\linewidth}p{0.32\linewidth}p{0.1\linewidth}p{0.16\linewidth}}
\toprule
Metric & Formula & Value & Reference \\
\midrule
Statistical Parity (\texttt{SP\_diff}) & $P(\hat Y{=}1|G{=}0)-P(\hat Y{=}1|G{=}1)$ & $-0.365$ & Dwork et al.\ (2012) \\
Disparate Impact (\texttt{DI\_ratio}) & $P(\hat Y{=}1|G{=}0)/P(\hat Y{=}1|G{=}1)$ & $0.408$ & Feldman et al.\ (2015) \\
Equal Opp., TPR (\texttt{EqOp\_plus\_diff}) & $\mathrm{TPR}(0)-\mathrm{TPR}(1)$ & $-0.296$ & Hardt et al.\ (2016) \\
Equal Opp., TNR (\texttt{EqOp\_minus\_diff}) & $\mathrm{TNR}(0)-\mathrm{TNR}(1)$ & $0.298$ & Hardt et al.\ (2016) \\
Eq.\ Odds, TPR (\texttt{EqOdds\_plus\_diff}) & $\mathrm{TPR}(0)-\mathrm{TPR}(1)$ & $-0.296$ & Hardt et al.\ (2016) \\
Eq.\ Odds, FPR (\texttt{EqOdds\_minus\_diff}) & $\mathrm{FPR}(0)-\mathrm{FPR}(1)$ & $-0.298$ & Hardt et al.\ (2016) \\
Treatment Equality, diff & $(\mathrm{FPR}/\mathrm{FNR})(0)-(\mathrm{FPR}/\mathrm{FNR})(1)$ & $-1.536$ & Berk et al.\ (2021) \\
Treatment Equality, ratio & $(\mathrm{FPR}/\mathrm{FNR})(0)/(\mathrm{FPR}/\mathrm{FNR})(1)$ & $0.140$ & Berk et al.\ (2021) \\
Overall Acc., diff & $\mathrm{Acc}(0)-\mathrm{Acc}(1)$ & $0.047$ & Berk et al.\ (2021) \\
Overall Acc., ratio & $\mathrm{Acc}(0)/\mathrm{Acc}(1)$ & $1.070$ & Berk et al.\ (2021) \\
Equalized Disincentive & $[\mathrm{TPR}{-}\mathrm{FPR}](0)-[\mathrm{TPR}{-}\mathrm{FPR}](1)$ & $0.003$ & Youden (1950) \\
PPV parity & $\mathrm{PPV}(0)-\mathrm{PPV}(1)$ & $-0.039$ & Chouldechova (2017) \\
NPV parity & $\mathrm{NPV}(0)-\mathrm{NPV}(1)$ & $0.100$ & Chouldechova (2017) \\
\bottomrule
\end{tabular}
\caption{The 13 binary fairness metrics, with the value each takes on the mock loan-approval example (eps not applicable). Full citations in Section~\ref{sec:refs}.}
\end{table}

\subsection{How to read the worked example}

The algorithm approves loans for group $G=0$ far less often than for $G=1$
(Statistical Parity $=-0.365$; Disparate Impact $=0.408$, well below the
conventional 0.8 ``four-fifths rule'' threshold). Part of this gap reflects
a real difference in creditworthiness across groups (the true repayment rate
is 0.356 vs.\ 0.564), but Equal Opportunity ($-0.296$) shows that even among
applicants who \emph{would} repay, $G=0$ is approved substantially less often
-- evidence of disparate treatment beyond what the underlying risk alone
would justify. Equalized Odds fails on both components (TPR and FPR), so
Equal Opportunity alone is not being satisfied ``by accident''.

Interestingly, Overall Accuracy is nearly identical across groups (in fact
slightly \emph{higher} for $G=0$, diff $=0.047$) and Equalized Disincentive
(the group-level Youden's J gap) is essentially zero ($0.003$). This is the
guide's central methodological point: \textbf{a single metric is not enough}.
Accuracy parity, or a near-zero Youden's-J gap, can coexist with large gaps
in Statistical Parity, Equal Opportunity, Disparate Impact, PPV and NPV --
each metric encodes a different normative notion of what ``fair'' means, and
they are not implied by one another (Kleinberg, Mullainathan \& Raghavan,
2016, formalize exactly this kind of incompatibility).

Finally, PPV parity ($-0.039$) and NPV parity ($0.100$) show that among
\emph{approved} applicants, those from $G=0$ are somewhat less likely to
actually repay than approved $G=1$ applicants -- the approval decision is
less ``reliable'' when applied to $G=0$.

\subsection{Conditional Equal Opportunity}

Conditional Equal Opportunity re-computes the Equal Opportunity TPR gap
separately within each level of a third variable (here, \texttt{age\_band}),
which is useful to check whether an aggregate gap is driven by one stratum or
is pervasive across all of them. This is an extension implemented in this
toolkit, building on Hardt et al.\ (2016); it is not itself a metric with a
single canonical source paper.

\begin{center}
\begin{tabular}{lccc}
\toprule
Age band & $\mathrm{TPR}(G{=}0)$ & $\mathrm{TPR}(G{=}1)$ & CEO gap \\
\midrule
18--30 & 0.445 & 0.758 & $-0.313$ \\
31--50 & 0.474 & 0.753 & $-0.279$ \\
51+    & 0.457 & 0.765 & $-0.307$ \\
\bottomrule
\end{tabular}
\end{center}

The gap is essentially constant across strata ($\approx -0.30$ everywhere),
so the aggregate Equal Opportunity violation is pervasive, not driven by a
single age group.

\section{The 4 multidimensional fairness metrics}

These extend fairness beyond two groups to an arbitrary number of subgroups
\texttt{ZG}. All four are governed by a discrimination-aversion parameter
$\varepsilon$; the values below use $\varepsilon = 0.10$ on the four
\texttt{ZG} subgroups of the mock data (\texttt{ZG} $=2\cdot G+\texttt{region}$).

\begin{table}[h]
\centering
\small
\begin{tabular}{p{0.22\linewidth}p{0.4\linewidth}p{0.08\linewidth}p{0.06\linewidth}p{0.16\linewidth}}
\toprule
Metric & Formula & Value & Fair? & Reference \\
\midrule
Statistical Subgroup Parity (\texttt{SubPar}) &
  $\max_m\big[P(zg_m)\cdot|P(\hat Y{=}1)-P(\hat Y{=}1|zg_m)|-\varepsilon\big]$ &
  $0.049$ & Yes & Kearns et al.\ (2018) \\
False Positive Subgroup Parity (\texttt{FalPos}) &
  $\max_m\big[P(Y{=}0,zg_m)\cdot|\mathrm{FPR}_{\text{all}}-\mathrm{FPR}_{zg_m}|-\varepsilon\big]$ &
  $0.021$ & Yes & Kearns et al.\ (2018) \\
Differential Fairness, $c=1$ (\texttt{DifFair}) &
  $\max_m P(\hat Y{=}c|zg_m)/\min_m P(\hat Y{=}c|zg_m)-e^{\varepsilon}$ &
  $2.606$ & No & Foulds et al.\ (2020) \\
Differential Fairness, $c=0$ &
  (same formula, $c=0$) &
  $2.044$ & No & Foulds et al.\ (2020) \\
Worst-Case Fairness, $c=1$ (\texttt{WorCas}) &
  $1-\min_m P(\hat Y{=}c|zg_m)/\max_m P(\hat Y{=}c|zg_m)$ &
  $0.616$ & No & Ghosh et al.\ (2021) \\
Worst-Case Fairness, $c=0$ &
  (same formula, $c=0$) &
  $0.511$ & No & Ghosh et al.\ (2021) \\
\bottomrule
\end{tabular}
\caption{The 4 multidimensional fairness metrics on the mock data ($\varepsilon = 0.10$). ``Fair?'' reports whether the metric's own fairness condition holds at that $\varepsilon$.}
\end{table}

Note the disagreement between the two ``subgroup'' metrics (SubPar, FalPos --
both pass at $\varepsilon=0.10$) and the two ``ratio''-based metrics
(DifFair, WorCas -- both fail). SubPar and FalPos weight each subgroup's
deviation by its population share, so a disparity concentrated in a smaller
subgroup contributes less to the global statistic; DifFair and WorCas instead
look at the most extreme ratio across \emph{all} subgroups regardless of
size, and are far more sensitive to exactly this kind of concentrated
disparity (the two largest \texttt{ZG} subgroups here, 2 and 3, have approval
rates around 0.60--0.63 vs.\ around 0.24--0.26 for subgroups 0 and 1). This is
the same lesson as Section 3.1 at a different level: which metric you use
determines what kind of unfairness you can detect.

\section{References}
\label{sec:refs}

\begin{itemize}
\item Berk, R., Heidari, H., Jabbari, S., Kearns, M., \& Roth, A.\ (2021).
  Fairness in Criminal Justice Risk Assessments: The State of the Art.
  \emph{Sociological Methods \& Research}, 50(1), 3--44.
\item Chouldechova, A.\ (2017). Fair Prediction with Disparate Impact: A
  Study of Bias in Recidivism Prediction Instruments. \emph{Big Data}, 5(2),
  153--163.
\item Dwork, C., Hardt, M., Pitassi, T., Reingold, O., \& Zemel, R.\ (2012).
  Fairness Through Awareness. \emph{Proceedings of the 3rd Innovations in
  Theoretical Computer Science Conference (ITCS)}.
\item Feldman, M., Friedler, S. A., Moeller, J., Scheidegger, C., \&
  Venkatasubramanian, S.\ (2015). Certifying and Removing Disparate Impact.
  \emph{Proceedings of the 21st ACM SIGKDD International Conference on
  Knowledge Discovery and Data Mining (KDD)}, 259--268.
\item Foulds, J. R., Islam, R., Keya, K. N., \& Pan, S.\ (2020). An
  Intersectional Definition of Fairness. \emph{2020 IEEE 36th International
  Conference on Data Engineering (ICDE)}.
\item Ghosh, A., Genuit, L., \& Reagan, M.\ (2021). Characterizing
  Intersectional Group Fairness with Worst-Case Comparisons. \emph{Proceedings
  of the 2nd Workshop on Diversity in Artificial Intelligence (DAI)}.
\item Hardt, M., Price, E., \& Srebro, N.\ (2016). Equality of Opportunity in
  Supervised Learning. \emph{Advances in Neural Information Processing
  Systems (NeurIPS)}.
\item Kearns, M., Neel, S., Roth, A., \& Wu, Z. S.\ (2018). Preventing
  Fairness Gerrymandering: Auditing and Learning for Subgroup Fairness.
  \emph{Proceedings of the 35th International Conference on Machine Learning
  (ICML)}. Companion empirical study: Kearns, Neel, Roth \& Wu, An Empirical
  Study of Rich Subgroup Fairness for Machine Learning, \emph{FAT* 2019}.
\item Kleinberg, J., Mullainathan, S., \& Raghavan, M.\ (2016). Inherent
  Trade-Offs in the Fair Determination of Risk Scores. \emph{arXiv:1609.05807}.
\item Youden, W. J.\ (1950). Index for Rating Diagnostic Tests. \emph{Cancer},
  3(1), 32--35.
\end{itemize}

\section*{Acknowledgments}

Alexandros Puente Pomar (Research Assistant) developed the original
implementation of these metrics that this toolkit builds on. His work -- the
correctness and care in the underlying computations -- was excellent, and
this public repository would not exist without it.

\section*{How to cite this toolkit}

Salas-Rojo, P.\ (2026). \emph{Fairness Metrics Toolkit}
[Software]. \url{https://github.com/pedrosalasrojo/fairness_metrics}

\end{document}
